\chapter{Identifying Parenting Values}

\begin{quote}
\itshape ``What matters is not just that our children act in certain ways, but that they have a commitment to certain values.''\\
\upshape --- Alfie Kohn
\end{quote}

\newpage
\section{What Are Parenting Values?}

Parenting is often described as a journey without a map, filled with daily decisions that shape the lives of our children.
In the absence of a definitive guide, our parenting values act as a compass, providing direction and purpose amidst the chaos \citep{kohn2005}.
Unlike specific parenting goals---such as getting a child to sleep through the night or to eat their vegetables---values are enduring beliefs about what is fundamentally important in life \citep{schwartz2012}.
They represent our deepest aspirations for the kind of relationship we want to cultivate with our children and the qualities we hope to model for them.

\subsection{The Origins of Our Values}

To fully understand what could lead to values, we must examine both our innate psychological needs and our developmental history.
According to Self-Determination Theory, all humans have basic psychological needs for autonomy, competence, and relatedness \citep{ryan2017}.
Our core values are often born out of our intrinsic drive to fulfill these essential needs within the context of our early relationships.
Furthermore, our family of origin plays a pivotal role; the values our caregivers modeled, whether explicitly taught or implicitly conveyed through their actions, deeply influence the moral template we develop \citep{hoffman2000}.
Values are not arbitrary rules; they are the crystallized expressions of what we have learned is necessary for survival, connection, and flourishing.

From the perspective of Nonviolent Communication (NVC), values are intimately connected to our universal human needs \citep{rosenberg2003}.
When we identify a parenting value, such as ``authenticity'' or ``compassion,'' we are essentially identifying the core needs that we are striving to meet in our family life.
Research in psychology underscores the importance of value-driven behavior.
Parents who have a clear understanding of their values are more resilient in the face of stress and are less likely to react impulsively or punitively when challenged by their children's behavior \citep{brown2012}.
By anchoring our parenting in deeply held values, we shift our focus from immediate compliance to long-term relational connection.

Furthermore, children learn values not primarily through lectures, but through observation.
Neurological research on the mirror-neuron system suggests that children are biologically wired to mimic the emotional states and actions of their primary caregivers \citep{rizzolatti2004}.
When parents consistently embody their values---such as demonstrating empathy during a conflict rather than resorting to punishment---they provide a powerful, experiential template for their children's developing moral framework \citep{hoffman2000}.

\begin{keyinsight}
Parenting values are enduring beliefs that serve as an internal compass, guiding our decisions and shaping our family culture.
By anchoring our actions in core values rather than short-term goals, we foster a deep, resilient connection with our children and model the qualities we hope to instill.
\end{keyinsight}

\section{Values Versus Expectations}

A critical step in identifying true parenting values is distinguishing them from societal expectations or internalized cultural pressures.
Many parents unknowingly adopt a set of ``shoulds'' based on how they were raised or what they perceive as normal in their community \citep{siegel2003}.
These externally imposed expectations often masquerade as values.

\subsection{Culture and Values}

It is impossible to separate our parenting values from the cultural waters in which we swim.
Different cultures emphasize distinctly different values; for instance, some prioritize interdependence and collective harmony, while others highly value individualism and personal achievement \citep{schwartz2012}.
While honoring one's cultural heritage is deeply meaningful, parents must be careful not to blindly accept cultural norms that conflict with their intrinsic need for empathetic connection.
NVC encourages us to examine cultural expectations critically, asking whether a specific cultural norm serves life and fosters mutual respect, or whether it relies on coercive power \citep{rosenberg2003}.
By doing so, parents can consciously choose which cultural values to pass down and which to thoughtfully leave behind.

For instance, a parent might believe they value ``obedience,'' but upon deeper reflection, they may realize this focus stems from a fear of societal judgment rather than a genuine desire for their child's authentic development.

When parenting is driven by external expectations, it inherently becomes conditional.
The parent's approval is unconsciously tied to the child's ability to meet specific, often rigid standards \citep{assor2004}.
This conditional regard undermines the child's intrinsic motivation and can lead to heightened anxiety and a fragile sense of self-worth \citep{ryan2017}.
In contrast, intrinsic values are self-chosen and reflect the parent's authentic self.
When we parent from a place of intrinsic values, such as valuing ``mutual respect'' or ``curiosity,'' we are able to offer our children unconditional positive regard, which is essential for healthy emotional development \citep{rogers1961}.

NVC teaches us to translate these external ``shoulds'' into language that connects with universal human needs \citep{rosenberg2003}.
Instead of operating from the expectation that a child ``should be quiet,'' a parent might identify a personal need for ``peace and order.''
By shifting the focus from the child's behavior to the parent's underlying need, the parent can approach the situation with empathy rather than blame.
This process of unearthing intrinsic values requires a courageous willingness to challenge inherited scripts and prioritize genuine connection over societal approval.

\begin{keyinsight}
It is essential to distinguish authentic parenting values from external societal expectations.
While expectations often lead to conditional parenting and anxiety, self-chosen intrinsic values foster unconditional positive regard and support the child's healthy, authentic development.
\end{keyinsight}

\section{Uncovering Your Core Values}

Identifying your core parenting values is an active, ongoing process of self-reflection and mindful awareness.
It requires parents to cultivate a non-judgmental curiosity about their own emotional triggers and deepest longings \citep{kabatzinn1990}.
One effective method for uncovering these values is to examine moments of intense emotional resonance, both positive and negative.
When we feel profound joy or deep satisfaction in our parenting, it is often an indicator that our core values are being honored.
Conversely, moments of intense frustration, guilt, or anger frequently signal that a deeply held value or need has been violated \citep{ekman1992}.

Another powerful tool for identifying values is reflecting on our own family of origin.
By thoughtfully exploring how we were parented, we can identify which values we wish to carry forward and which generational patterns we consciously choose to disrupt \citep{siegel2003}.
This process of making the implicit explicit allows parents to heal from past relational trauma and create a more intentional, empathetic family culture \citep{vanderkolk2014}.
Furthermore, engaging in mindful self-reflection practices, such as journaling or mirror meditation, can significantly enhance a parent's capacity for objective self-awareness \citep{well2022}.

Ultimately, uncovering core values is not about creating a rigid set of rules, but about developing a flexible, living framework that guides our interactions.
As children grow and family dynamics evolve, our understanding of our values will inevitably deepen and shift.
By maintaining an open, empathetic dialogue with ourselves and our children, we can ensure that our parenting remains anchored in the timeless principles of connection, compassion, and love \citep{rosenberg2003}.

\begin{keyinsight}
Uncovering core values is an ongoing process of self-reflection, requiring mindful awareness of our emotional triggers and inherited family patterns.
By actively exploring these areas, parents can develop an intentional, flexible framework that prioritizes deep relational connection.
\end{keyinsight}

\section{Judgments and Values}

A profound challenge in identifying and living by our parenting values is distinguishing between a genuine value and a moralistic judgment.
In NVC, \citet{rosenberg2003} differentiates between "value judgments," which reflect our beliefs about how life can best be served, and "moralistic judgments," which imply wrongness or badness on the part of people who do not act in harmony with our values.
When we operate from moralistic judgments, we use labels like "lazy," "selfish," or "disrespectful" to describe a child's behavior.
This language of blame creates defensiveness and disconnection, making it nearly impossible to impart the very values we hold dear.

Instead, NVC suggests expressing our values through the lens of our own feelings and unmet needs.
If a parent values cooperation but a child refuses to clean up, a moralistic judgment would be, "You are so irresponsible."
A value-based expression, however, focuses on the parent's experience: "When I see the toys on the floor, I feel frustrated because I need support in keeping our shared space tidy."
By separating our core values from punitive judgments, we maintain an empathetic connection, increasing the likelihood that our children will voluntarily embrace the values we are attempting to model \citep{kohn2005}.

\begin{keyinsight}
Moralistic judgments assign blame and create disconnection, whereas value judgments reflect our beliefs about what serves life.
By translating our judgments into expressions of our own feelings and needs, we can model our values empathetically without resorting to blame or coercion.
\end{keyinsight}

\newpage
\section{Practice Exercise}
\begin{practiceexercise}
\textbf{Multiple Choice Questions}
\begin{enumerate}
  \item What is the primary function of parenting values?
    \begin{itemize}
      \item[A)] To establish strict rules for children to follow daily.
      \item[B)] To guarantee that children meet developmental milestones.
      \item[C)] To act as an internal compass guiding relational decisions.
      \item[D)] To ensure parents receive approval from their community.
    \end{itemize}
  \vspace{\fill}
  \item How do children primarily learn values according to neurological research?
    \begin{itemize}
      \item[A)] Through memorization of household rules and chores.
      \item[B)] Through experiential observation and mirroring of caregivers.
      \item[C)] Through repeated verbal instructions and lectures.
      \item[D)] Through a system of carefully planned rewards and punishments.
    \end{itemize}
  \vspace{\fill}
  \item From an NVC perspective, what are values closely related to?
    \begin{itemize}
      \item[A)] Societal expectations and cultural norms.
      \item[B)] Behavioral goals and specific parenting techniques.
      \item[C)] Universal human needs.
      \item[D)] Arbitrary judgments of right and wrong behavior.
    \end{itemize}
\end{enumerate}
\vspace{\fill}

\newpage

\textbf{Explanations for Questions 1--3}

\textbf{Question 1: C.} Parenting values serve as a guiding compass, providing direction and purpose in the complex daily interactions with children, focusing on long-term relationship building rather than short-term behavioral compliance.
\vspace{\fill}

\textbf{Question 2: B.} Neurological research on the mirror-neuron system demonstrates that children are biologically wired to mimic and absorb the emotional states and actions they observe in their primary caregivers.
\vspace{\fill}

\textbf{Question 3: C.} In Nonviolent Communication, the values we hold dear are ultimately expressions of our deep, universal human needs, such as the need for connection, authenticity, and peace.
\vspace{\fill}
\newpage

\begin{enumerate}
  \setcounter{enumi}{3}
  \item What distinguishes a genuine parenting value from a societal expectation?
    \begin{itemize}
      \item[A)] Genuine values are easier to enforce consistently.
      \item[B)] Genuine values are intrinsically chosen rather than externally imposed.
      \item[C)] Societal expectations always prioritize the child's feelings.
      \item[D)] Societal expectations are based on universal human needs.
    \end{itemize}
  \vspace{\fill}
  \item What is a common consequence of parenting driven by external expectations?
    \begin{itemize}
      \item[A)] The development of unconditional positive regard.
      \item[B)] Conditional parenting that undermines the child's intrinsic motivation.
      \item[C)] A profound sense of deep connection and mutual respect.
      \item[D)] The elimination of all behavioral challenges in the home.
    \end{itemize}
  \vspace{\fill}
  \item How can NVC help parents address internalized cultural ``shoulds''?
    \begin{itemize}
      \item[A)] By helping parents enforce the ``shoulds'' more compassionately.
      \item[B)] By teaching parents to ignore all community standards.
      \item[C)] By translating expectations into language about underlying needs.
      \item[D)] By proving that all cultural norms are inherently harmful.
    \end{itemize}
\end{enumerate}
\vspace{\fill}

\newpage
\textbf{Explanations for Questions 4--6}

\textbf{Question 4: B.} Authentic values reflect a parent's intrinsic, self-chosen beliefs about what matters most, whereas societal expectations are often adopted unconsciously to avoid judgment or conform to external norms.
\vspace{\fill}

\textbf{Question 5: B.} When parents prioritize external standards, their approval often becomes tied to the child's compliance, resulting in conditional regard that can increase anxiety and damage a child's sense of self-worth.
\vspace{\fill}

\textbf{Question 6: C.} NVC encourages parents to dig beneath rigid rules and expectations to uncover the universal human needs driving those desires, shifting the focus from controlling the child to empathetic understanding.
\vspace{\fill}
\newpage

\begin{enumerate}
  \setcounter{enumi}{6}
  \item What can intense emotional frustration in parenting frequently indicate?
    \begin{itemize}
      \item[A)] That the parent is naturally unsuited for caregiving.
      \item[B)] That the child is intentionally acting with malice.
      \item[C)] That a deeply held personal value or need has been violated.
      \item[D)] That the parent needs to establish stricter punishments.
    \end{itemize}
  \vspace{\fill}
  \item Why is reflecting on one's family of origin important for identifying values?
    \begin{itemize}
      \item[A)] It allows parents to replicate their childhood experiences exactly.
      \item[B)] It provides a rigid template for resolving all family conflicts.
      \item[C)] It helps parents consciously choose which patterns to keep or disrupt.
      \item[D)] It proves that all parenting challenges are the fault of the grandparents.
    \end{itemize}
  \vspace{\fill}
  \item How should a parent's understanding of their core values be viewed?
    \begin{itemize}
      \item[A)] As a permanent, unchanging list established before birth.
      \item[B)] As a rigid set of rules that must never be compromised.
      \item[C)] As a flexible, living framework that evolves with the family.
      \item[D)] As a private code that should never be discussed with children.
    \end{itemize}
\end{enumerate}
\vspace{\fill}

\newpage
\textbf{Explanations for Questions 7--9}

\textbf{Question 7: C.} Strong emotional reactions, such as profound frustration or anger, act as vital signals that our underlying needs are not being met or that our core values are somehow being compromised in the moment.
\vspace{\fill}

\textbf{Question 8: C.} By exploring how they were parented, individuals can make implicit generational patterns explicit, empowering them to intentionally adopt positive values and heal from past traumas.
\vspace{\fill}

\textbf{Question 9: C.} Because children develop and family dynamics constantly shift, our understanding of our values must remain a dynamic, flexible guide rather than a stagnant set of unyielding rules.
\vspace{\fill}
\end{practiceexercise}

\newpage
\section{Reflection Questions}
\begin{reflection}
\small
\begin{enumerate}[itemsep=0pt, parsep=0pt, topsep=0pt]
  \item What are three core values you hope your children observe in your daily actions?
  \item Can you identify a parenting ``should'' you follow because of external pressure rather than internal belief?
  \item Think of a recent moment of intense parenting joy; what underlying value was being honored?
  \item When you feel deeply frustrated with your child, which of your universal needs is usually unmet?
  \item How do your current parenting values differ from the values emphasized in your family of origin?
  \item In what situations do you find yourself prioritizing immediate behavioral compliance over long-term relational connection?
  \item How might you translate a specific expectation you have for your child into a statement of your own needs?
  \item What is one mindful practice you could adopt to become more aware of your emotional triggers?
  \item How does viewing values as a ``flexible framework'' change your approach to making difficult parenting decisions?
\end{enumerate}
\end{reflection}